% =====================================================================
%  IEEE Journal of Photovoltaics -- Supplementary Material
%  Single-column, since supplements are not typeset in the journal's
%  two-column body. Same class so fonts and sizes match the manuscript.
%  BUILD:  ./build.sh supp
% =====================================================================
\documentclass[lettersize,journal,onecolumn]{IEEEtran}

\usepackage{amsmath,amsfonts,amssymb}
\usepackage{array}
\usepackage{textcomp}
\usepackage{url}
% let long repository URLs break rather than run into the margin
\def\UrlBreaks{\do\/\do-\do.\do:\do_}
\Urlmuskip=0mu plus 1mu\relax
\usepackage{graphicx}
\usepackage{cite}
\usepackage{xcolor}
\usepackage{booktabs}
\usepackage[margin=1in]{geometry}

\newif\ifedits
\editstrue          % <-- \editsfalse strips every editorial note
\newcommand{\edit}[1]{\ifedits{\color{red}\textbf{[#1]}}\fi}

\newcommand{\FASnI}{FASnI\textsubscript{3}}
\newcommand{\CABI}{Cu\textsubscript{2}AgBiI\textsubscript{6}}
\newcommand{\BZS}{BaZrS\textsubscript{3}}
\newcommand{\CABB}{Cs\textsubscript{2}AgBiBr\textsubscript{6}}
\newcommand{\TiOx}{TiO\textsubscript{2}}
\newcommand{\percm}{\,mA\,cm\textsuperscript{$-$2}}

\renewcommand{\thefigure}{S\arabic{figure}}
\renewcommand{\thetable}{S\arabic{table}}
\renewcommand{\thesection}{S\arabic{section}}

\begin{document}

\title{Supplementary Material\\
\large Interface Band Offsets, Not Optical Collection, Limit Four Chemically
Distinct Lead-Free Absorbers}
\author{Md.~Mohiuddin and Alamgir~Kabir}
\maketitle

\section{Supplementary Figures}

\begin{figure}[!ht]
\centering
\includegraphics[width=\textwidth]{figures_si/figS1_optical_replacement.pdf}
\caption{\textbf{Optical-constant source selection.} Extinction coefficient
$k$ of (a)~\FASnI{} and (b)~\BZS, comparing candidate datasets. (a)~A
pixel-digitized reconstruction (rejected) splices two traced figures
of~\cite{ghimire2017} at $1.70\,$eV; across that $5\,\mathrm{nm}$ join $k$
falls $0.4263\rightarrow0.0792$, a $5.8\times$ discontinuity, after which it
sits nearly flat to $850\,$nm -- behaviour no direct-gap absorber exhibits.
The adopted dataset is the tabulated dielectric function from the same
paper's supplementary information (695 points, point-by-point inversion of
measured ellipsometry). (b)~A Kramers--Kronig reconstruction of a
\emph{computed} $\varepsilon_2$ (rejected) is hard-zeroed at $690\,$nm,
cutting from $\alpha = 2.6\times10^{4}\,\mathrm{cm^{-1}}$ directly to zero
inside its own absorption range. The adopted dataset is the measured
ellipsometric dielectric function~\cite{nishigaki2020}.}
\end{figure}

\begin{figure}[!ht]
\centering
\includegraphics[width=\textwidth]{figures_si/figS2_validation.pdf}
\caption{\textbf{Validation.} (a)~Per-wavelength energy-budget residual for all
four absorbers against the $\pm2\,\%$ acceptance band; the worst value across
244 solves is $-1.43\,\%$, at $300\,$nm, the finest mesh in the study.
(b)~Band-seam agreement: wavelengths at the $400$ and $600\,$nm band
boundaries solved on both neighbouring meshes, which are generated
independently. Agreement is $0.020$--$0.096\,\%$. This is the check that
detects the Floquet periodic-mesh defect of Section~S4; the energy residual
does not.}
\end{figure}

\begin{figure}[!ht]
\centering
\includegraphics[width=\textwidth]{figures_si/figS3_subgap_cu2agbii6.pdf}
\caption{\textbf{\CABI{} sub-gap absorption.} (a)~Absorber absorptance with the
$2.06\,$eV Tauc gap marked; the shaded region is sub-gap. (b)~Decomposition of
the $15.083$\percm{} optical ceiling by spectral region; the segments
partition the total exactly. $61.7\,\%$ lies above the Tauc gap; $29.0\,\%$
is an Urbach-like tail with characteristic energy $\approx88\,$meV;
$3.3\,\%$ a weak tail; and $6.0\,\%$ ($0.908$\percm) arises
where $\alpha$ stops decaying and sits flat at
$\approx2.7\times10^{3}\,\mathrm{cm^{-1}}$ for $130\,$nm before stepping to
zero at $886\,$nm, a construction artifact of the optical file rather than a
physical tail.}
\end{figure}

\begin{figure}[!ht]
\centering
\includegraphics[width=\textwidth]{figures_si/figS4_voc_grid_cu2agbii6.pdf}
\caption{\textbf{\CABI{} open-circuit voltage is not grid-converged.} (a)~$J$--$V$
on a mixed $1$--$25\,\mathrm{mV}$ continuation grid versus a uniform
$1\,\mathrm{mV}$ grid. The curves are identical through the maximum power
point ($P_\mathrm{max}=10.761\,\mathrm{mW\,cm^{-2}}$,
$V_\mathrm{mpp}=0.786\,$V, $J_\mathrm{mpp}=13.691$\percm, to five significant
figures). (b)~Magnified pre-$V_\mathrm{oc}$ tail: the curves diverge only
beyond $\approx1.3\,$V where $J<0.2$\percm{} and the characteristic is nearly
flat, so a small numerical difference moves the zero crossing from $1.515\,$V
to $1.384\,$V.}
\end{figure}

\begin{figure}[!ht]
\centering
\includegraphics[width=\textwidth]{figures_si/figS5_brad_bazrs3.pdf}
\caption{\textbf{\BZS{} radiative coefficient.} (a)~van Roosbroeck--Shockley
$B_\mathrm{rad}$ versus the lower integration limit, for the rejected
(computed) and adopted (measured) optical constants. A candidate value of
$2.335\times10^{-9}\,\mathrm{cm^3\,s^{-1}}$ corresponds to the rejected
optics integrated from $1.75\,$eV while $n_i^2$ is built from
$E_g=1.88\,$eV -- internally inconsistent, and $\approx12\times$ too large
even on its own data. Evaluated consistently at $1.88\,$eV on the measured
optics, $B_\mathrm{rad}=5.37\times10^{-12}\,\mathrm{cm^3\,s^{-1}}$
(adopted). (b)~Device $J$--$V$ with both coefficients: a $435\times$ change
moves $V_\mathrm{oc}$ by $3\,$mV, establishing that \BZS's $V_\mathrm{oc}$
is not radiatively limited.}
\end{figure}

\begin{figure}[!ht]
\centering
\includegraphics[width=\textwidth]{figures_si/figS6_generation_map.pdf}
\caption{\textbf{Depth- and wavelength-resolved generation.} Absorbed power
density $Q_h(z,\lambda)$ through the $300\,$nm absorber under unit incident
field, plotted from the illuminated face. Each absorber's band edge is
directly visible as the wavelength beyond which generation ceases:
$\approx880\,$nm for \FASnI, $\approx700\,$nm for \CABI{} with its sub-gap
tail continuing weakly to $900\,$nm, $\approx660\,$nm for \BZS, and
$\approx570\,$nm for \CABB. The diagonal fringes are Fabry--P\'erot
interference in the finite-thickness stack, which a Beer--Lambert treatment
cannot reproduce. The contrast between panels explains the collection
behaviour reported in the main text: \FASnI{} generates strongly across the
full depth, so an optical improvement that front-loads generation moves
carriers into a region it collects less efficiently, whereas \CABB{} generates
only in a narrow band near the illuminated face regardless.}
\end{figure}

\begin{figure}[!ht]
\centering
\includegraphics[width=\textwidth]{figures_si/figS7_tmm_crosscheck.pdf}
\caption{\textbf{Transfer-matrix verification of the electromagnetic solver.}
Absorber absorptance from the finite-element solve (open circles) against an
independent scalar transfer-matrix calculation with identical optical inputs
and identical geometry (solid lines), for all four absorbers. Worst-case
disagreement across all $4\times61$ solves is $7\times10^{-4}$ in
absorptance and $0.02\,\%$ in AM1.5G photocurrent. For \CABB{} the dashed
curve shows the thick-contact (semi-infinite Au) limit: the visible departure
in the $430$--$570\,$nm window is the back-contact truncation artifact of
Section~S3, which the identical-geometry agreement isolates as geometric
rather than numerical.}
\end{figure}

\clearpage
\section{Device architecture and parameter provenance}

The stack is fixed for all four devices: glass / FTO ($600\,$nm) / \TiOx{}
($50\,$nm) / absorber ($300\,$nm) / Spiro-OMeTAD ($200\,$nm) / Au. The optical
unit cell is $350\times350\,\mathrm{nm}^2$ with Floquet periodicity on all
lateral faces.

\textbf{Sources.} Absorber optics: \FASnI~\cite{ghimire2017},
\CABI~\cite{kamppinen2026,sansom2021}, \BZS~\cite{nishigaki2020},
\CABB~\cite{eddekkar2024}. Gaps and masses: \BZS~\cite{kumar2021}, band
alignment~\cite{eya_dzade}; \CABB~\cite{eddekkar2024,filip2016}. Lifetimes:
\BZS{} $1$\,ns from single-crystal transient photoluminescence
\cite{nielsen2025} and $33$\,ns computed from the sulfur-interstitial
capture coefficients of~\cite{yuan2024}. \FASnI{} electron affinity:
experimental photoelectron spectroscopy in air (PESA)~\cite{wu2022chi}, with
the $4.0$\,eV sweep endpoint from~\cite{scaps_fasni3}. Transport layers:
\cite{rahman2023,usman2024}. Contact optics: FTO~\cite{vonrottkay1996},
\TiOx~\cite{devore1951}, Au~\cite{mcpeak2015}.

\textbf{Provenance codes.} \textbf{[EXP]} experimental; \textbf{[THEO]}
computed, level stated; \textbf{[SIM-LIT]} from published device simulations;
\textbf{[DERIVED]} computed here from literature inputs by a stated
parameter-free relation; \textbf{[SWEPT]} treated as a sensitivity axis.

\begin{table}[!ht]
\caption{Absorber parameters}
\centering\small
\begin{tabular}{lllllc}
\toprule
Quantity & \BZS & \FASnI & \CABB & \CABI & Status \\
\midrule
$E_g$ (eV)        & 1.88 & 1.41 & 2.32 & 2.03$^{a}$ & [THEO]/[EXP] \\
$\chi$ (eV)       & 4.1 (3.9, 4.0) & 3.50 (3.75, 4.0) & 4.0 (3.79, 4.19) & 3.22 (3.16, 3.36) & \textbf{[SWEPT]}$^{b}$ \\
$m_e^*/m_0$       & 0.3  & 0.3  & 0.33 & 0.6  & [THEO] \\
$m_h^*/m_0$       & 0.9  & 0.9  & 0.84 & 1.0  & [THEO] \\
$\tau_\mathrm{SRH}$ & 1/33\,ns & 43/123\,ns & 13.7\,ns/1\,$\mu$s & 33/3.3\,ns & \textbf{[SWEPT]} \\
$\varepsilon_r$   & 79.4 & 8.2 (6, 10) & 5.8 (3.54) & 6.3 (12) & \textbf{[SWEPT]}$^{b}$ \\
$\mu_n$, $\mu_p$ (cm$^2$V$^{-1}$s$^{-1}$) & 37, 11 & 67 (22) & 0.57 (6) & 1.7 (5.1) & \textbf{[SWEPT]}$^{b}$ \\
Doping (cm$^{-3}$) & $N_d$ $4.4{\times}10^{10}$ & $N_a$ $1.5{\times}10^{17}$ & $N_a$ $1.42{\times}10^{12}$ & $N_a$ $1{\times}10^{12}$ & [SIM-LIT] \\
$B_\mathrm{rad}$ (cm$^3$s$^{-1}$) & $\mathbf{5.37{\times}10^{-12}}$ & $7.99{\times}10^{-8}$ & $3.19{\times}10^{-12}$ & $1.51{\times}10^{-10}$ & \textbf{[DERIVED]} \\
\bottomrule
\end{tabular}

\vspace{2pt}
{\footnotesize All values read back from the saved models. $^{b}$Format:
base value (alternates swept). The near-intrinsic \BZS{} donor density is
the deliberate literature value for optimized low-defect processing; the
conventional high-temperature-sulfurization range ($10^{19}$--$10^{20}\,
\mathrm{cm^{-3}}$) exceeds $N_c$ and is incompatible with the Boltzmann
statistics used here.}

\vspace{2pt}
{\footnotesize $^{a}$Transport-model value, read back from the saved model;
the 2.06\,eV Tauc gap of~\cite{sansom2021} is used for the sub-gap
decomposition of Fig.~S3.}
\end{table}

\begin{table}[!ht]
\caption{Transport-layer and contact parameters}
\centering\small
\begin{tabular}{lllc}
\toprule
Quantity & \TiOx{} (ETL) & Spiro-OMeTAD (HTL) & Source \\
\midrule
$E_g$ (eV)                        & 3.2  & 2.9  & \cite{usman2024,rahman2023} \\
$\chi$ (eV)                      & 4.15 & 2.15 & \cite{usman2024,rahman2023} \\
$\varepsilon_r$                  & --   & 3    & \cite{rahman2023} \\
$\mu_n$, $\mu_p$ (cm$^2$V$^{-1}$s$^{-1}$) & -- & $1\times10^{-4}$ & \cite{rahman2023} \\
$N_c$, $N_v$ (cm$^{-3}$)          & $2.2\times10^{18}$, $1\times10^{19}$ & $2.2\times10^{18}$, $2.2\times10^{18}$ & \cite{rahman2023,usman2024} \\
Doping (cm$^{-3}$)                & $N_d$ $5\times10^{17}$ & $N_a$ $1\times10^{18}$ & \cite{rahman2023} \\
Optical                           & Sellmeier, $k=0$~\cite{devore1951} & $n=1.75$, $k=0$ & Sec.~S6 \\
FTO work function (eV)            & \multicolumn{2}{l}{4.6} & \cite{rahman2023} \\
\bottomrule
\end{tabular}
\end{table}

\textbf{Disclosed disagreements.} Two literature conflicts are carried openly
rather than silently resolved. Spiro-OMeTAD band gap: $2.9$\,eV
\cite{rahman2023} versus $3.2$\,eV \cite{tiwari2022}; Spiro $N_v$:
$2.2\times10^{18}$ \cite{rahman2023} versus $1.8\times10^{19}$
\cite{tiwari2022}. We use \cite{rahman2023} throughout. A third reported
value, a Spiro hole mobility of $2\times10^{4}\,\mathrm{cm^2V^{-1}s^{-1}}$
\cite{tiwari2022}, is rejected as a transcription error: it is roughly ten
orders of magnitude above the well-established
${\sim}10^{-4}\,\mathrm{cm^2V^{-1}s^{-1}}$ for this material. Both
disagreements and the rejection were confirmed against each paper's own
tables. The HTL acceptor doping is the saved-model
value ($1\times10^{18}\,\mathrm{cm^{-3}}$); see the reproducibility caution
below.

\textbf{A reproducibility caution.} The HTL acceptor doping is
$1\times10^{18}\,\mathrm{cm^{-3}}$ in the saved model but
$1\times10^{19}\,\mathrm{cm^{-3}}$ in the build script that nominally
creates it; a later stage changed the value and only the model file records
it. The tenfold difference eases hole extraction enough to shift \BZS's
fill factor from $0.34$ to $0.38$, so parameters must be read back from the
saved model, never from the build script.

\subsection{Gap consistency across the three roles a gap plays}
A band gap is used in three distinct roles here (the transport model's
$n_i^2$, the detailed-balance ceiling, and the absorption onset of the optical
constants) and these are not automatically the same number. Measured directly
from the computed spectra:

\begin{table}[!ht]
\caption{Band gap used in each role, and the resulting spread}
\centering\small
\begin{tabular}{lrrrr}
\toprule
Absorber & transport $n_i^2$ & DB ceiling & optical onset & spread \\
\midrule
\FASnI & 1.41 & 1.41 & 1.39 & 0.02\,eV \\
\BZS   & 1.88 & 1.94 & 1.88 & 0.06\,eV \\
\CABB  & 2.32 & 2.10 & 2.21 & \textbf{0.22\,eV} \\
\CABI  & 2.03 & 2.06 & 1.41 & \textbf{0.65\,eV} \\
\bottomrule
\end{tabular}
\end{table}

\FASnI{} and \BZS{} are consistent. \textbf{\CABB{} is not}: its HSE06
transport gap ($2.32\,$eV), the $2.10\,$eV used for its detailed-balance
ceiling, and the $2.21\,$eV onset of its own optical constants disagree by up
to $0.22\,$eV. The device consequence was tested directly by re-solving the device with
$E_g = 2.21\,$eV (its own optical onset): $J_\mathrm{sc}$, $P_\mathrm{max}$
and PCE are unchanged to four decimal places ($1.170\,\%$), while the gap
shift appears entirely in $V_\mathrm{oc}$ ($-111\,$mV $\approx \Delta E_g$,
as $n_i^2$ dictates) with the fill factor compensating exactly
($0.281\rightarrow0.300$) because the maximum power point sits on the
barrier-limited plateau at $0.60\,$V. The same test on \BZS{}
($1.88\rightarrow1.94\,$eV) gives the same structure: PCE unchanged at
$4.499\,\%$, $V_\mathrm{oc}$ $+62\,$mV, FF $0.340\rightarrow0.325$.
Efficiency-level conclusions are therefore invariant to the gap choice;
only $V_\mathrm{oc}$ carries it, by a predictable ${\approx}\Delta E_g$. \textbf{\CABI's $0.65\,$eV spread is by construction}, since its
optical file carries a deliberate sub-gap tail to $1.40\,$eV; the consequence
there is material and is discussed in the main text.



\section{The back-contact truncation artifact in \CABB}
The optical model truncates the Au contact at $45\,$nm backed by a perfect
electric conductor -- a deliberate cost measure, justified by Au's
$12$--$17\,$nm near-infrared skin depth (three-plus skin depths of metal,
below $10^{-6}$ returning power). The justification fails in one window: Au's
interband damping dips near $450$--$570\,$nm, roughly doubling the skin
depth, and \CABB{} is the one absorber nearly transparent exactly there.
Light reaches the back contact, and the PEC backing returns an unphysically
strong wave: at $510\,$nm the truncated geometry gives $R = 0.469$ while
realistically thick Au gives $R = 0.292$. The other three absorbers
extinguish this window before it reaches the contact and agree with the
thick-contact limit to $\leq0.35\,\%$ in photocurrent.

The correction uses the transfer-matrix method of Fig.~S7, which is exact for
this planar stack and validated against the finite-element solver to
$7\times10^{-4}$ on identical geometry. Three gates were passed in order:
(i) the PEC-geometry TMM reproduces the FEM spectrum to $9\times10^{-4}$
across absorptance and reflectance jointly ($6.5\times10^{-4}$ on absorber
absorptance alone -- the metric of Table~S5 and Fig.~S7); (ii) the same pipeline reproduces the FEM generation
profile $G(z)$, including the AM1.5G bin weighting, to $0.10\,\%$ with
identical photocurrent; (iii) only then was the thick-contact spectrum and
$G(z)$ emitted and the \CABB{} transport re-solved on it. Corrected numbers:
ceiling $2.546\rightarrow2.440$\percm{} (the truncated value overstated
the ceiling by $4.3\,\%$), device
$J_\mathrm{sc}$ $2.473\rightarrow2.371$\percm, PCE
$1.22\rightarrow1.17\,\%$; collection is unchanged ($97.2\,\%$). The
$45\,$nm-truncation figure of $2.546$ appears nowhere in the main text
except as the disclosed pre-correction value.

\section{The Floquet periodic-mesh defect}
Floquet periodic boundary conditions require the paired faces of the unit cell
to carry \textbf{identical surface meshes}. A free-tetrahedral mesher does not
produce this. The solver does not raise an error; it interpolates across the
mismatch, and the computed per-layer absorption depends on the particular mesh
realization. Measured on this stack at $700\,$nm:

\begin{table}[!ht]
\caption{Absorber absorptance versus mesh realization, before the fix}
\centering\small
\begin{tabular}{rr}
\toprule
Elements & $A_\mathrm{absorber}$ \\
\midrule
120\,153 & 0.719 \\
120\,346 & 0.663 \\
123\,031 & 0.463 \\
\bottomrule
\end{tabular}
\end{table}

A $36\,\%$ swing for a $2.4\,\%$ change in element count. The remedy is to
mesh $x=0$ with a free-triangular operation, copy that mesh onto $x=L$, repeat
for $y$, and only then fill with tetrahedra; feature order matters because the
$x$ and $y$ faces share the cell's vertical edges. After the fix the same
comparison gives $0.783103$ versus $0.782888$ ($0.03\,\%$). The global
energy budget also improves, from $1.2\,\%$ to $0.03\,\%$ -- but note that
even the defective mesh sat comfortably inside the $\pm2\,\%$ acceptance
band while the layer partition was wrong by $36\,\%$: the residual shifts,
but far too weakly to be diagnostic, which is precisely why it cannot serve
as the detector.

\textbf{The energy-conservation residual does not detect this defect.} It
compares total absorbed power against total incident flux and is insensitive
to how that total is partitioned between layers, which is precisely what the
photocurrent depends on. Detection requires either a mesh-versus-mesh
comparison or a tangential-field equality check on the paired faces (which
reaches $0.00000\,\%$ after the fix). Two traps arise in the latter: compare
the \emph{tangential} field components, not $|E|$, because the normal
component sampled on a face is interpolated from whichever side the evaluator
picks; and expect the $x$ and $y$ pairs to differ, because the incident wave
is $x$-polarized.

\section{Numerical procedure}
\textbf{Wavelength banding.} Element count scales as $\lambda^{-3}$, so meshes
are generated per band: $300$--$390$, $400$--$590$, $600$--$900\,$nm. Element
size is $\lambda/(n_\mathrm{eff}\,\mathrm{ppw})$ with
$n_\mathrm{eff}=|\tilde{n}|$, subject to a skin-depth clamp and a floor of
three elements across any layer thickness. The wavelength-independent cap is
applied only to air and the perfectly matched layer, never to the absorber.

\textbf{Two independent photocurrent routes.} The depth integral of $G(z)$ and
the spectral integral of layer absorptance are computed separately and
compared every run. They share the solve but nothing else. Agreement across
the four absorbers is $+0.05$, $+0.03$, $+0.04$, $-0.04\,\%$.

\textbf{AM1.5G quadrature.} Irradiance is integrated over each wavelength bin
rather than sampled at bin centres, and the outer bin edges are clamped to the
grid endpoints. Bin-centre sampling biased the result by $\approx10\,\%$ and
unclamped edges by a further $+23\,\%$ on a coarse grid; both are quadrature
errors rather than physics, and both were exposed by the two-route comparison.

\textbf{Bias continuation.} For \CABI{} a $25\,$mV grid converges onto spurious
S-shaped negative-differential-resistance branches reporting $\mathrm{FF}=4.13$
and $\mathrm{PCE}=86.6\,\%$; a $1\,$mV grid removes them. Only dense results
are quoted for that absorber. The artifact is reproducible rather than
stochastic.

\begin{table}[!ht]
\caption{Validation ledger}
\centering\small
\begin{tabular}{lrrrr}
\toprule
Check & \FASnI & \CABI & \BZS & \CABB \\
\midrule
Optical budget closure       & $-0.00\,\%$ & $+0.00\,\%$ & $+0.00\,\%$ & $+0.00\,\%$ \\
Worst energy residual        & $-1.36\,\%$ & $-1.39\,\%$ & $-1.43\,\%$ & $-1.35\,\%$ \\
$G(z)$ vs spectrum           & $+0.05\,\%$ & $+0.03\,\%$ & $+0.04\,\%$ & $-0.04\,\%$ \\
Band seam, $400\,$nm         & $0.070\,\%$ & $0.079\,\%$ & $0.075\,\%$ & $0.096\,\%$ \\
Band seam, $600\,$nm         & $0.023\,\%$ & $0.020\,\%$ & $0.020\,\%$ & $0.000\,\%$ \\
TMM vs FEM, max $|\Delta A|$ (identical geometry) & $7.0\times10^{-4}$ & $6.2\times10^{-4}$ & $5.9\times10^{-4}$ & $6.5\times10^{-4}$ \\
\bottomrule
\end{tabular}
\end{table}

\section{Optical model limitations}
\textbf{Lossless transport layers.} \TiOx{} is modeled with a Sellmeier real
index (measured on rutile; ETL films are typically anatase, a further
approximation of this treatment) clamped below $400\,$nm (its $283\,$nm pole otherwise returns $n=5.58$
at the blue end of the grid), and Spiro-OMeTAD with a constant $n=1.75$; both
carry $k=0$. Their budget entries are therefore exactly zero and all front
parasitic absorption is attributed to the FTO. \TiOx{} genuinely absorbs below
its $\approx390\,$nm edge, so UV generation in the absorber is slightly
overestimated. The treatment is identical for all four devices, so relative
comparisons are unaffected.

\textbf{FTO.} Previously modeled with a fabricated flat $k=0.02$; here it uses
measured SnO$_2$:F dispersion.

\textbf{Gold contact.} Truncated to $45\,$nm in the optical model, more than
three skin depths at all wavelengths considered and backed by a perfect
electric conductor, while the transport model retains the full $80\,$nm. The
generation profile is therefore shifted by $35\,$nm when handed to the
transport stage, and both frames are provided in the archive.

\section{The Cu$_2$AgBiI$_6$ sub-gap collection bracket}
$38.2\,\%$ of \CABI's photocurrent originates below its Tauc gap (Fig.~S3),
and the transport model collects it like band-to-band generation. To bound
what that assumption is worth, the device was re-solved on two restricted
generation profiles built from the per-wavelength absorbed-power data (the
full reconstruction reproduces the production profile exactly before any
restriction is applied): one with only above-gap bins
($\lambda \leq 600\,$nm), and one with all bins except the $750$--$880\,$nm
segment where the absorption coefficient sits at an unphysical flat floor --
a construction artifact of the optical file rather than a real Urbach tail
(these are the $10\,$nm solve bins covering the flat segment, whose true edge
in the optical file lies at ${\approx}885\,$nm; the removed bins carry
$0.99$\percm{} versus the $0.91$\percm{} continuum integral of Fig.~S3, the
difference being bin quantization). The same quantization applies to the
above-gap member: its binned profile ($\lambda \leq 600$\,nm bins)
integrates to $9.484$\percm{} versus the $9.31$\percm{} continuum integral
of Fig.~S3, and the device's $9.442$\percm{} is $99.6\,\%$ of its own
binned profile, not of the continuum figure -- no entry exceeds its
ceiling.

\begin{table}[!ht]
\caption{\CABI{} sub-gap collection bracket (uniform 1\,mV grid)}
\centering\small
\begin{tabular}{lrrrr}
\toprule
Generation & $J_\mathrm{sc}$ & $V_\mathrm{oc}$ (V) & FF & PCE \\
\midrule
above-gap only (floor)        &  9.442 & 1.367 & 0.541 & $6.98\,\%$ \\
artifact-free (defensible)    & 14.032 & 1.382 & 0.521 & $10.10\,\%$ \\
full collection (as Table~I, main text) & 15.017 & 1.384 & 0.518 & $10.76\,\%$ \\
\bottomrule
\end{tabular}
\end{table}

The decomposition separates two very different quantities: the construction
artifact alone is worth $0.66$ absolute PCE points, and the assumption that
genuine tail states collect like band carriers a further $3.12$ points. The
artifact-free $10.10\,\%$ is the defensible central value; Table~I of the
main text retains
the full-collection number for cross-absorber consistency, with this bracket
qualifying it.

\section{Interface recombination is not representable in this formulation}
No experimental interface-recombination velocity exists for any of these
absorber/transport-layer pairs, and the efficiencies in the main text
correspond to defect-free interfaces; they are upper bounds.

An attempt was made to sweep an interface recombination velocity
$S = 10^2$--$10^6\,\mathrm{cm\,s^{-1}}$ at the two junctions, and it is
reported here as unrepresentable rather than as a result: the solver's
trap-assisted surface-recombination boundary feature, offered every boundary
in the geometry, accepts only exterior and lateral boundaries and excludes
the two interior heterojunction boundaries entirely, silently solving with an
empty applicable selection if assigned to them. Two consequences are drawn.
(i) All efficiencies here are explicitly labelled defect-free upper bounds,
with no interface-sensitivity claim attached. (ii) Because the limiting
mechanism in every device is carrier accumulation at a band-offset cliff,
real interface recombination would act on precisely the populations this
model shows piling up at those interfaces; the upper-bound label is
therefore load-bearing, not a formality. Representing junction recombination
faithfully would require a different construction (for example, thin
defective interlayers), which we leave to future work.

\section{Data and code availability}
All inputs, scripts, results and figure code are archived at\\
{\footnotesize\url{https://github.com/mohiuddinsciencelover-glitch/leadfree-pv-3d-simulation}}
\texttt{REPRODUCE.md} documents the pipeline end to end, including the five
traps that cost significant time during development: the periodic-mesh defect,
the build-script/model parameter divergence, the \CABI{} continuation-grid
artifact, the two AM1.5G quadrature errors, and the Au-truncation artifact of
Section~S3.

\bibliographystyle{IEEEtran}
\bibliography{references}

\end{document}
